% ========================================================================
% TEMPLATE LATEX PARA TCC NO FORMTATO DE ARTIGO
% Baseado na classe abntex2 com personalizações para o IFPI
% Conforme e Manual de Normalização de Trabalhos Acadêmicos do IFPI - 2024 
% e as normas ABNT NBR:
%   14724:2011
%   6024:2012
%   6027:2012
%   10520:2023
%   6023:2018 
%   15287:2011 
% ========================================================================

% ========================================================================
% PREÂMBULO E CONFIGURAÇÕES
% ========================================================================

% ========================================================================
% CLASSE DO DOCUMENTO
% ========================================================================
\documentclass[
    12pt,        % Tamanho da fonte principal
    openany,     % Capitulos podem iniciar em qualquer pagina
    oneside,     % Impressao apenas frente
    a4paper,     % Formato A4
    english,     % Idioma adicional
    french,      % Idioma adicional
    spanish,     % Idioma adicional
    brazil,      % Idioma principal
    article      % Formato de artigo
]{abntex2}

% ========================================================================
% PACOTES FUNDAMENTAIS
% ========================================================================

% Codificacao e fontes
\usepackage[utf8]{inputenc}
\usepackage[T1]{fontenc}
\usepackage{lmodern}
\usepackage{times}

% Idioma
\usepackage[brazil]{babel}

% Regras ABNT / referencias
\usepackage[alf]{abntex2cite}
\usepackage{config/abntex-ifpi}

% ========================================================================
% PACOTES PARA FORMATACAO E LAYOUT
% ========================================================================
\usepackage{microtype}
\usepackage{indentfirst}
\usepackage{color}
\usepackage{graphicx}
\usepackage{float}
\usepackage{amsmath,amsfonts,amssymb}
\usepackage{booktabs}
\usepackage{array}
\usepackage{longtable}
\usepackage{multirow}
\usepackage{rotating}
% \usepackage{siunitx} % alinhamento numérico com separador decimal em vírgula
% \sisetup{
%     output-decimal-marker = {,},
%     input-decimal-markers = {,},
%     group-minimum-digits = 3,
%     group-separator = {.},
%     table-number-alignment = right
% }
% \newcolumntype{d}[1]{S[table-format=#1, input-symbols = {-}]}

% Listas, siglas e simbolos
\usepackage{enumitem}
\usepackage{acronym}
\usepackage{nomencl}

% URLs
\usepackage{url}

% ========================================================================
% MARGENS, ESPACAMENTO E NUMERACAO
% ========================================================================
\setlrmarginsandblock{3cm}{2cm}{*}
\setulmarginsandblock{3cm}{2cm}{*}
\checkandfixthelayout

\SingleSpacing
\setlength{\parindent}{1.25cm}

\setcounter{secnumdepth}{3}
\setcounter{tocdepth}{3}

% ========================================================================
% LISTAS E NOMES PADRAO
% ========================================================================
\renewcommand{\listfigurename}{LISTA DE ILUSTRACOES}
\renewcommand{\listtablename}{LISTA DE TABELAS}
\renewcommand{\nomname}{LISTA DE ABREVIATURAS E SIGLAS}

% ========================================================================
% LEGENDAS E AUTOREF
% ========================================================================
\renewcommand{\figureautorefname}{Figura}
\renewcommand{\tableautorefname}{Tabela}
\renewcommand{\equationautorefname}{Equacao}
\renewcommand{\quadroautorefname}{Quadro}
% ========================================================================
% COMANDOS UTILITARIOS
% ========================================================================
\newcommand{\citep}[2]{\cite[p.~#2]{#1}}
\newcommand{\citeonlinep}[2]{\citeonline[p.~#2]{#1}}
\newcommand{\citemulti}[1]{\cite{#1}}


\newcommand{\fontetabela}[1]{%
    % \vspace{0.2cm}
    \begin{center}
        {\footnotesize Fonte: #1}
    \end{center}
}

% ========================================================================
% HYPERREF (adiado para depois dos dados do trabalho)
% ========================================================================
\AtBeginDocument{%
    \hypersetup{%
        pdftitle={\imprimirtitulo},
        pdfauthor={\imprimirautor},
        pdfsubject={\imprimirpreambulo},
        pdfcreator={LaTeX with abnTeX2},
        pdfkeywords={abnt}{latex}{abntex}{abntex2}{projeto de pesquisa}{IFPI},
        colorlinks=false,
        hidelinks=true,
        linkcolor=black,
        citecolor=black,
        filecolor=black,
        urlcolor=black,
    }%
}

% ========================================================================
% AJUSTES FINAIS
% ========================================================================
\sloppy
\hyphenpenalty=5000
\tolerance=1000

\setlist{nosep}
\renewcommand{\thepage}{\arabic{page}}
   %configurações de formatação

% ========================================================================
% DADOS DO ARTIGO CIENTÍFICO
% Edite os valores abaixo para personalizar o template.
% ========================================================================

% Instituicao e identificacao do curso
\instituicao{Instituto Federal de Educação, Ciência e Tecnologia do Piauí}
\campus{Campus Floriano}
\curso{Tecnologia em Análise e Desenvolvimento de Sistemas}
\areaconcentracao{} % Opcional: defina a area de concentracao ou deixe vazio

% Titulo e autoria
\titulo{titulo}
\subtitulo{}
\tituloestrangeiro{} % Título traduzido (ex: em inglês para o abstract)
\autor{autor}
\emailautor{autor@aluno.ifpi.edu.br}
\vinculoautor{IFPI -- Campus Floriano}


% Orientação
\orientador{Prof. Me. Ronaldo Pires Borges}
\emailorientador{ronaldo.borges@ifpi.edu.br}
\vinculoorientador{IFPI -- Campus Floriano} % Vínculo do orientador

% Coorientação (deixe vazio se não houver coorientador)
\coorientador{}
\emailcoorientador{}
\vinculocoorientador{}


% Informacoes gerais
\local{Floriano -- PI}
\data{2026}
% \data{\the\year} % auto
\dataaprovacao{XX/XX/XXXX} % ESCREVA A DATA DA BANCA

% Descrição apresentada na folha de rosto
\preambulo{Trabalho de Conclusão de Curso (Artigo Científico) apresentado como exigência parcial para obtenção do diploma do Curso de Tecnologia em Análise e Desenvolvimento de Sistemas do Instituto Federal de Educação, Ciência e Tecnologia do Piauí, Campus Floriano.}
% Natureza do trabalho 
% \naturezatrabalho{Projeto de pesquisa apresentado ao Instituto Federal de Educacao, Ciencia e Tecnologia do Piaui como requisito para aprovação na disciplina...}

% Grau pretendido (opcional)
% \grau{} 
% \grau{Tecnólogo em Analise e Desenvolvimento de Sistemas} 

% ========================================================================
% INSTRUCOES RAPIDAS
% - Substitua os textos em caixa alta pelos dados reais.
% - Campos vazios podem ser removidos ou comentados.
% - Para atualizar o resumo e demais secoes, edite os arquivos em principal.tex.
% ========================================================================
           %dados institucionais e do projeto

% ========================================================================
% INÍCIO DO DOCUMENTO
% ========================================================================

\begin{document}
% Seleciona o idioma do documento (conforme pacotes do babel)
\selectlanguage{brazil}
% Retira espaço extra obsoleto entre as frases
\frenchspacing

% ========================================================================
% ELEMENTOS PRÉ-TEXTUAIS
% ========================================================================
% ========================================================================
% ELEMENTOS PRÉ-TEXTUAIS
% Capa e folha de rosto conforme artigo TCC IFPI
% ========================================================================

% ========================================================================
% CAPA (Obrigatória em artigos do IFPI segundo o manual TCC)
% ========================================================================
\imprimircapa

% ========================================================================
% FOLHA DE ROSTO (Obrigatória)
% ========================================================================
\imprimirfolhaderosto*

% ========================================================================
% AGRADECIMENTOS (OPCIONAL - EM ARTIGOS DEVE IR NO FINAL)
% ========================================================================
% O manual do IFPI diz que agradecimentos em artigos vão como elemento pós-textual,
% portanto mude-os para o final do `main.tex` se for usar.

% ========================================================================
% CABEÇALHO DO ARTIGO (Título, autores, e-mails)
% ========================================================================
\imprimircabecalhoartigo

% ========================================================================
% RESUMO EM PORTUGUÊS E PALAVRAS-CHAVE
% ========================================================================
\setlength{\absparsep}{18pt} % Ajusta espaçamento entre parágrafos do resumo

\begin{resumo}
    Este artigo científico tem como objetivo [DESCREVER O OBJETIVO PRINCIPAL DA PESQUISA]. A metodologia utilizada baseia-se em [DESCREVER A METODOLOGIA: pesquisa bibliográfica, estudo de caso, pesquisa experimental, etc.]. Como referencial teórico, serão utilizados os trabalhos de [CITAR PRINCIPAIS AUTORES DA ÁREA]. A pesquisa justifica-se pela [APRESENTAR A JUSTIFICATIVA E RELEVÂNCIA DO TEMA]. Os resultados incluem [DESCREVER OS RESULTADOS ALCANÇADOS]. Este estudo contribui para [APRESENTAR A CONTRIBUIÇÃO PARA A ÁREA DE CONHECIMENTO].
    
    \vspace{\onelineskip}
    \noindent
    \textbf{Palavras-chave}: palavra-chave 1; palavra-chave 2; palavra-chave 3; palavra-chave 4; palavra-chave 5.
\end{resumo}

% ========================================================================
% RESUMO EM INGLÊS (ABSTRACT) E KEYWORDS
% ========================================================================
\begin{resumo}[ABSTRACT]
    \begin{otherlanguage*}{english}
        \setlength{\parindent}{0pt}% Garante que não haverá recuo na primeira linha do abstract
        This scientific article aims to [DESCRIBE THE MAIN OBJECTIVE OF THE RESEARCH]. The methodology used is based on [DESCRIBE THE METHODOLOGY: bibliographic research, case study, experimental research, etc.]. As theoretical framework, the works of [CITE MAIN AUTHORS IN THE FIELD] will be used. The research is justified by [PRESENT THE JUSTIFICATION AND RELEVANCE OF THE THEME]. Results include [DESCRIBE REACHED RESULTS]. This study contributes to [PRESENT THE CONTRIBUTION TO THE KNOWLEDGE AREA].
        
        \vspace{\onelineskip}
        \noindent
        \textbf{Keywords}: keyword 1; keyword 2; keyword 3; keyword 4; keyword 5.
    \end{otherlanguage*}
\end{resumo}

% ========================================================================
% DATA DE APROVAÇÃO
% ========================================================================
\vspace{\onelineskip}
\noindent
Data de aprovação: \@dataaprovacao

\vspace{1cm}

% ========================================================================
% CONFIGURAÇÕES PARA INÍCIO DO TEXTO
% ========================================================================
% Impede quebra de página padrão do abntex2 após pré-textuais
\makeatletter
\renewcommand{\textual}{%
  \pagestyle{ifpihead}
  \aliaspagestyle{chapter}{ifpihead}
  \nouppercaseheads
  \bookmarksetup{startatroot}
}
\makeatother

\textual


% ========================================================================
% ELEMENTOS TEXTUAIS
% ========================================================================

\section{INTRODUÇÃO}
\label{sec:introducao}

A introdução de um projeto de pesquisa deve apresentar o tema de forma clara e contextualizada, demonstrando sua relevância no cenário atual. Este capítulo tem como objetivo situar o leitor no contexto da pesquisa, apresentando o problema a ser investigado e justificando a importância do estudo.

O tema desta pesquisa insere-se no campo de [ÁREA DE CONHECIMENTO], especificamente abordando [TEMA ESPECÍFICO]. A escolha deste tema justifica-se pela [JUSTIFICATIVA INICIAL], considerando sua relevância tanto no âmbito acadêmico quanto na aplicação prática.

Segundo \citeonline{marconi2017}, a área de [ÁREA DE CONHECIMENTO] tem apresentado crescente interesse por estudos que abordem [ASPECTO RELEVANTE DO TEMA]. Esta tendência evidencia a necessidade de pesquisas que contribuam para o avanço do conhecimento nesta área.

O presente projeto de pesquisa propõe investigar [OBJETO DE ESTUDO], com foco em [ASPECTO ESPECÍFICO A SER INVESTIGADO]. A pesquisa será desenvolvida considerando [CONTEXTO OU DELIMITAÇÃO DO ESTUDO], o que permitirá uma análise aprofundada do fenômeno em questão.


\section{PROBLEMA DE PESQUISA}
\label{sec:problema}

O problema de pesquisa constitui o elemento central que orienta todo o desenvolvimento do estudo. Conforme \citeonline{gil2019}, um problema bem formulado deve ser claro, preciso e operacional, permitindo a definição de objetivos específicos e a escolha de métodos adequados para sua investigação.

No contexto desta pesquisa, observa-se que [DESCRIÇÃO DO PROBLEMA IDENTIFICADO]. Esta situação tem gerado [CONSEQUÊNCIAS OU IMPACTOS DO PROBLEMA], evidenciando a necessidade de estudos que possam contribuir para sua compreensão e possível solução.

Diante do exposto, formula-se o seguinte problema de pesquisa:

\textbf{[PERGUNTA DE PESQUISA PRINCIPAL]?}

Este questionamento surge da observação de que [CONTEXTUALIZAÇÃO DO PROBLEMA], conforme evidenciado por \citeonline{silva2023} em seus estudos sobre [ÁREA RELACIONADA].


\section{HIPÓTESES}
\label{sec:hipoteses}   

As hipóteses representam respostas provisórias ao problema de pesquisa, baseadas no conhecimento teórico disponível e na experiência do pesquisador. Segundo \citeonline{marconi2017}, as hipóteses devem ser testáveis e formuladas de maneira clara e precisa.

Para o problema apresentado, formulam-se as seguintes hipóteses:


\textbf{Hipótese Principal (H1):} [HIPÓTESE PRINCIPAL DA PESQUISA]

\textbf{Hipóteses Secundárias:}
\begin{itemize}
    \item \textbf{H2:} [SEGUNDA HIPÓTESE]
    \item \textbf{H3:} [TERCEIRA HIPÓTESE, SE APLICÁVEL]
\end{itemize}

Essas hipóteses serão testadas através da metodologia proposta, permitindo confirmar ou refutar as suposições iniciais da pesquisa.



\section{OBJETIVOS}
\label{sec:objetivos}

Os objetivos definem o que se pretende alcançar com a realização da pesquisa. Devem ser formulados de forma clara e precisa, utilizando verbos no infinitivo que expressem ações concretas e mensuráveis.

\subsection{Objetivo Geral}
\label{subsec:objetivo-geral}

[DESCREVER O OBJETIVO GERAL DA PESQUISA, UTILIZANDO VERBO NO INFINITIVO E EXPRESSANDO DE FORMA CLARA O QUE SE PRETENDE ALCANÇAR COM O ESTUDO]

\subsection{Objetivos Específicos}
\label{subsec:objetivos-especificos}

\begin{enumerate}
    \item Objetivo específico 1;
    \item Objetivo específico 2;
    \item Objetivo específico 3;
    \item Objetivo específico 4 (se aplicável).
\end{enumerate}

\section{JUSTIFICATIVA}
\label{sec:justificativa}

A justificativa deve demonstrar a relevância e a importância da pesquisa, apresentando argumentos que evidenciem sua contribuição para o avanço do conhecimento científico e/ou para a solução de problemas práticos.

Esta pesquisa justifica-se por diversos aspectos:

\textbf{Relevância Científica:} O tema abordado contribui para o avanço do conhecimento na área de [ÁREA DE CONHECIMENTO], especialmente no que se refere a [ASPECTO ESPECÍFICO]. Conforme destacado por \citeonline{inep2022}, existe uma lacuna na literatura sobre [LACUNA IDENTIFICADA], que esta pesquisa pretende contribuir para preencher.

\textbf{Relevância Social:} Os resultados desta pesquisa poderão beneficiar [PÚBLICO-ALVO OU SOCIEDADE EM GERAL], uma vez que [EXPLICAR COMO OS RESULTADOS PODEM SER APLICADOS]. Segundo dados de \cite{instituicao2023}, [APRESENTAR DADOS QUE EVIDENCIEM A RELEVÂNCIA SOCIAL].

\textbf{Relevância Acadêmica:} Este estudo contribuirá para a formação acadêmica do pesquisador e para o desenvolvimento de pesquisas futuras na área, fornecendo subsídios teóricos e metodológicos para estudos similares.

\textbf{Viabilidade:} A pesquisa é viável considerando os recursos disponíveis, o tempo previsto para sua execução e o acesso às fontes de informação necessárias.

\section{REFERENCIAL TEÓRICO}
\label{sec:referencial}

O referencial teórico constitui a base conceitual da pesquisa, apresentando as teorias, conceitos e estudos que fundamentam a investigação. Este capítulo deve demonstrar o conhecimento do pesquisador sobre o tema e situar a pesquisa no contexto do conhecimento científico existente.

\subsection{[TÍTULO DA PRIMEIRA SEÇÃO TEÓRICA]}
\label{subsec:fundamentacao1}

[DESENVOLVER A FUNDAMENTAÇÃO TEÓRICA PRINCIPAL DO TEMA, APRESENTANDO OS CONCEITOS CENTRAIS E AS TEORIAS QUE EMBASAM A PESQUISA]

Segundo \citeonline{creswell2018}, [APRESENTAR CONCEITO OU TEORIA FUNDAMENTAL]. Esta definição é complementada por \citeonline{gil2019}, que acrescenta [INFORMAÇÃO COMPLEMENTAR].

\begin{figure}[htbp]
    \centering
    \caption{Exemplo de figura no documento}
    \includegraphics[width=0.3\textwidth]{img/Logo-IFPI-Floriano-Horizontal.png}
    \fonte{IFPI - Instituto Federal do Piauí.}
    \label{fig:exemplo}
\end{figure}

Quando elementos visuais, como figuras, quadros ou tabelas, forem inseridos no texto, devem ser obrigatoriamente identificados e mencionados no corpo do trabalho, de modo a estabelecer uma relação clara entre a informação apresentada e sua discussão. Conforme orienta a \citeonline{ABNT_NBR14724_2011},  esses elementos devem ser citados no texto por seu número e título, garantindo uniformidade e rastreabilidade. Assim, pode-se fazer referência, por exemplo, à \autoref{fig:exemplo}, na qual se apresenta a logomarca horizontal do IFPI Campus Floriano.

% ou
% \inserirfigura{img/Logo-IFPI-Floriano-Horizontal.png}{0.3\textwidth}{Exemplo de figura no documento}{exemplo}

% ========================================================================
% CITAÇÃO DIRETA LONGA - EXEMPLO
% ========================================================================
% Citação direta com mais de 3 linhas deve ser destacada com recuo de 4 cm da margem esquerda,
% sem aspas, com fonte menor (tamanho 10) e espaçamento simples.
% A formatação automática do abntex2 já aplica estas regras quando se usa o ambiente citacao.

\begin{citacao}
Quando uma citação direta ultrapassa três linhas do texto original, ela deve ser apresentada em destaque, com recuo especial da margem esquerda, sem o uso de aspas, utilizando fonte menor que a do corpo do texto e espaçamento simples entre as linhas. Esta formatação diferenciada permite ao leitor identificar claramente que se trata de um trecho literalmente reproduzido de outra fonte, mantendo a integridade do pensimento do autor ciado \cite{ABNT_NBR10520_2023}.
\end{citacao}

\subsection{ESTUDOS RELACIONADOS}
\label{subsec:estudos-relacionados} 

Esta seção apresenta uma revisão dos principais estudos relacionados ao tema da pesquisa, evidenciando o estado da arte e identificando lacunas que justifiquem a realização do presente estudo.

A \textbf{tabela} é uma forma não discursiva de apresentação de informações em que o dado numérico é o elemento central, organizada de modo a facilitar comparações, análises e leitura rápida dos resultados. Segundo as normas da ABNT, as tabelas devem apresentar título centralizado, estrutura aberta (sem bordas laterais ou inferiores fechadas), fonte abaixo e devem ser citadas no texto conforme sua numeração \cite{ABNT_NBR14724_2011,ABNT_NBR15287_2011,ABNT_NBR6024_2012}. Lembre-se que toda tabela, figura ou quadro inseridos no seu trabalhos dever ser citatdo no texto, como por exemplo, a \autoref{tab:estudos}, que exemplifica a formatação adequada de uma tabela.

\begin{table}[htbp] %"table" é o ambiente padrão do LaTeX para tabelas.
    \centering
    \caption{Exemplo de tabela }
    \label{tab:estudos}
    \begin{tabular}{p{4cm}cc}
        \hline
        \textbf{Estudo} & \textbf{Amostra} & \textbf{Métrica principal} \\ %\\ indica quebra de linha
        \hline\hline %Linha dupla para separar o cabeçalho do conteúdo
        Costa e Rocha (2023) & 120 respondentes & Média de adoção = 4,1 \\
        Silva e Santos (2023) & 18 escolas analisadas & Taxa de conclusão = 87\% \\
        Souza (2021) & 32 entrevistas & Índice de concordância = 0,74 \\
        \hline
    \end{tabular}
    \fonte{Elaborado pelo autor a partir dos estudos revisados.}
\end{table}

% =============== OBSERVAÇÃO SOBRE POSICIONAMENTO DE TABELAS E QUADROS ==============
% No ambiente table (tabela) e quadro (quadro), o LaTeX permite especificar opções de posicionamento.

% \begin{table}[htbp] <- Essas letras são indicadores de posicionamento:

% h	here = aqui	“coloque exatamente aqui se possível”
% t	top = topo	“se não couber aqui, tente no topo da página”
% b	bottom = rodapé	“ou no fim da página”
% p	page of floats	“ou faça uma página só de tabelas/figuras”

% Assim, [htbp] significa:

% “Tente colocar aqui; se não der, tenta no topo; depois no rodapé; depois numa página separada.”
% É o mais usado porque dá flexibilidade, o que evita que o LaTeX “quebre” seu layout.
% ====================================================================================

O \textbf{quadro}, diferentemente da tabela, é utilizado para apresentar dados qualitativos, como definições, características, categorias ou descrições. Ele possui estrutura fechada em todos os lados e organiza informações textuais de maneira comparativa e clara, sendo recomendado para sínteses conceituais ou classificações dentro do projeto de pesquisa \cite{ABNT_NBR14724_2011,ABNT_NBR15287_2011,ABNT_NBR6024_2012}. O \autoref{qua:estudos} exemplifica a formatação adequada de um quadro.

\begin{quadro}[htbp]  %"quadro" é o ambiente personalizado do abntex2 para quadros
    \centering
    \caption{Exemplo de quadro}
    \label{qua:estudos}
    \begin{tabular}{|p{3cm}|p{4cm}|p{6cm}|} % os caracteres | indicam as bordas verticais
        \hline
        \textbf{Autor/Ano} & \textbf{Metodologia} & \textbf{Principais Resultados} \\
        \hline
        Autor A (2023) & Pesquisa quantitativa & Ênfase na mensuração de impacto tecnológico \\ \hline
        Autor B (2022) & Estudo de caso & Identifica barreiras institucionais para inovação \\ \hline
        Autor C (2021) & Pesquisa qualitativa & Discute competências desenvolvidas pelos participantes \\ \hline
    \end{tabular}
    \fonte{Elaborado pelo autor.}
\end{quadro}

Você pode usar a feramenta \textit{Tables Genetator}\footnote{https://www.tablesgenerator.com/} para criar tabelas e quadros.

\section{METODOLOGIA}
\label{sec:metodologia}

A metodologia descreve os procedimentos que serão adotados para a realização da pesquisa, incluindo a caracterização do estudo, os métodos de coleta e análise de dados, e os aspectos éticos envolvidos.

\subsection{CARACTERIZAÇÃO DA PESQUISA}
\label{subsec:caracterizacao}

Esta pesquisa pode ser caracterizada, quanto aos seus objetivos, como [EXPLORATÓRIA/DESCRITIVA/EXPLICATIVA], uma vez que [JUSTIFICAR A CLASSIFICAÇÃO]. Segundo \citeonline{gil2019}, pesquisas [TIPO DA PESQUISA] têm como objetivo [OBJETIVO DO TIPO DE PESQUISA].

Quanto à abordagem, trata-se de uma pesquisa [QUALITATIVA/QUANTITATIVA/MISTA], pois [JUSTIFICAR A ABORDAGEM]. \citeonline{creswell2018} destaca que a abordagem [TIPO DE ABORDAGEM] é adequada quando [SITUAÇÕES EM QUE A ABORDAGEM É ADEQUADA].

Em relação aos procedimentos técnicos, será realizada uma [PESQUISA BIBLIOGRÁFICA/ESTUDO DE CASO/PESQUISA DE CAMPO/ETC.], que, conforme \citeonline{marconi2017}, caracteriza-se por [CARACTERÍSTICAS DO PROCEDIMENTO ESCOLHIDO].

\begin{figure}[htbp]
    \centering
    \caption{Figura que ilustre a sua metodologia}
    \includegraphics[width=0.9\textwidth]{img/tema do tcc.png}
    \fonte{Adaptado de \citeonline{Moretti_ComoEscolherTemaTCC_2024}}
    \label{fig:exemplo2}
\end{figure}

\subsection{POPULAÇÃO E AMOSTRA}
\label{subsec:populacao-amostra}

A população desta pesquisa é constituída por [DEFINIR A POPULAÇÃO]. Para a definição da amostra, será utilizada a técnica de [TIPO DE AMOSTRAGEM], que, segundo \citeonline{hair2019}, é adequada quando [JUSTIFICAR A ESCOLHA DA TÉCNICA DE AMOSTRAGEM].

O tamanho da amostra foi calculado considerando [CRITÉRIOS PARA CÁLCULO DA AMOSTRA], resultando em [NÚMERO] participantes/elementos.

\subsection{COLETA DE DADOS}
\label{subsec:coleta-dados}

A coleta de dados será realizada através de [INSTRUMENTOS DE COLETA], que permitirão obter informações sobre [TIPO DE INFORMAÇÕES A SEREM COLETADAS].

Os instrumentos escolhidos foram:
\begin{itemize}
    \item \textbf{[INSTRUMENTO 1]:} [DESCRIÇÃO E JUSTIFICATIVA]
    \item \textbf{[INSTRUMENTO 2]:} [DESCRIÇÃO E JUSTIFICATIVA]
\end{itemize}

\subsection{ANÁLISE DOS DADOS}
\label{subsec:analise-dados}

Os dados coletados serão analisados utilizando [TÉCNICAS DE ANÁLISE]. Para os dados quantitativos, será utilizada [ESTATÍSTICA DESCRITIVA/INFERENCIAL], com auxílio do software [NOME DO SOFTWARE]. Os dados qualitativos serão analisados através de [TÉCNICA DE ANÁLISE QUALITATIVA].

\subsection{ASPECTOS ÉTICOS}
\label{subsec:aspectos-eticos}

Esta pesquisa será desenvolvida respeitando os princípios éticos estabelecidos pela Resolução nº 466/2012 do Conselho Nacional de Saúde. [SE APLICÁVEL: O projeto será submetido ao Comitê de Ética em Pesquisa antes do início da coleta de dados.]

Todos os participantes serão informados sobre os objetivos da pesquisa e assinarão o Termo de Consentimento Livre e Esclarecido (TCLE), garantindo sua participação voluntária e o direito de desistir a qualquer momento.



\section{CONSIDERAÇÕES FINAIS}
\label{sec:consideracoes-finais}

As considerações finais são a afirmação sintética da ideia central do trabalho e dos pormenores apresentados no texto.
Devem ser feitas com base nos objetivos propostos e devem ser coerentes com todo o corpo textual.

% Os itens de Cronograma e Recursos (geralmente presentes em Projetos de Pesquisa) 
% foram comentados ou embutidos como subseções/comentários para preservar a estrutura de Artigo Científico.
% Em Artigos (resultados finais), o foco recai sobre Discussão dos Resultados.



% ========================================================================
% ELEMENTOS PÓS-TEXTUAIS
% ========================================================================
% ========================================================================
% ELEMENTOS PÓS-TEXTUAIS
% Referências, apêndices e anexos conforme ABNT NBR 14724:2011
% ========================================================================

% ========================================================================
% REFERÊNCIAS BIBLIOGRÁFICAS
% ========================================================================
\postextual

% Configura o título das referências
\bibliography{referencias}

% ========================================================================
% GLOSSÁRIO (OPCIONAL)
% ========================================================================
% Descomente se necessário incluir glossário
% \glossary

% ========================================================================
% APÊNDICES (OPCIONAL)
% ========================================================================
% Descomente e personalize se houver apêndices
\quebrapagina
\begin{apendicesenv}
    % Imprime uma página indicando o início dos apêndices
    % \partapendices
    
    % ========================================================================
    % APÊNDICE A - EXEMPLO DE APÊNDICE
    % ========================================================================
    \chapter{TÍTULO DO PRIMEIRO APÊNDICE}
    \label{apendice:primeiro}
    
    Este é um exemplo de apêndice. Os apêndices são textos ou documentos elaborados pelo próprio autor, a fim de complementar sua argumentação, sem prejuízo da unidade nuclear do trabalho.
    
    Exemplos de conteúdo para apêndices:
    \begin{itemize}
        \item Questionários utilizados na pesquisa
        \item Roteiros de entrevistas
        \item Códigos de programas desenvolvidos
        \item Cálculos detalhados
        \item Tabelas extensas com dados brutos
        \item Figuras complementares
    \end{itemize}
    
    % ========================================================================
    % APÊNDICE B - EXEMPLO DE QUESTIONÁRIO
    % ========================================================================
    % \chapter{QUESTIONÁRIO APLICADO NA PESQUISA}
\end{apendicesenv}

% ========================================================================
% FIM DO DOCUMENTO
% ========================================================================
\end{document}
